%Version 3.1 December 2024
%%%%%%%%%%%%%%%%%%%%%%%%%%%%%%%%%%%%%%%%%%%%%%%%%%%%%%%%%%%%%%%%%%%%%%
%% PV056 ML & DM -- Plant Disease Classification (Topic 2)           %%
%% Triplet-loss metric learning + Mahalanobis unknown detection      %%
%%%%%%%%%%%%%%%%%%%%%%%%%%%%%%%%%%%%%%%%%%%%%%%%%%%%%%%%%%%%%%%%%%%%%%

\documentclass[pdflatex,sn-mathphys-num]{sn-jnl}

\usepackage{graphicx}
\usepackage{multirow}
\usepackage{amsmath,amssymb,amsfonts}
\usepackage[title]{appendix}
\usepackage{booktabs}

\raggedbottom

\begin{document}

\title[Plant Disease Classification with Triplet Loss]{Plant Disease Classification with Triplet-Loss Metric Learning and Mahalanobis Unknown-Class Detection}

\author*[1]{\fnm{Jozef} \sur{Kraus}}\email{jozefkraus@mail.muni.cz}
\author[1]{\fnm{David} \sur{Charv\'at}}\email{536483@mail.muni.cz}
\author[1]{\fnm{Richard} \sur{Schwarz}}\email{rschwarz@mail.muni.cz}
\author[1]{\fnm{Mat\v{e}j} \sur{Kub\'ik}}\email{514143@mail.muni.cz}

\affil*[1]{\orgdiv{Faculty of Informatics}, \orgname{Masaryk University}, \orgaddress{\city{Brno}, \country{Czech Republic}}}

\maketitle

\section{Introduction and Related Work}\label{sec:intro}

We address plant-disease image classification on the PlantVillage dataset \cite{mohanty2016using} under two regimes: closed-set recognition of 15 disease/healthy classes and open-set detection of a held-out ``unknown'' disease (\textit{Tomato Bacterial spot}). CNN classifiers on PlantVillage routinely exceed 99\% accuracy on its own test split \cite{ferentinos2018deep}; this work is therefore concerned less with squeezing out the remaining points of accuracy and more with two orthogonal questions: (i) how to learn a \emph{representation} that is useful beyond a single softmax head, and (ii) how to decide that a query does not belong to any known class. We tackle (i) via triplet-loss metric learning \cite{schroff2015facenet,hermans2017defense} with a prototype (nearest-centroid) classifier, and (ii) via class-conditional Mahalanobis distance in the learned embedding space \cite{lee2018simple}, an approach the task statement explicitly prefers over softmax-based open-set scores, which are known to be poorly calibrated for novel classes. We note one important caveat: PlantVillage images are captured in lab conditions on uniform backgrounds, and it has been shown that classifiers exploit this background shortcut \cite{noyan2022uncovering}, so high accuracies on PlantVillage should not be read as a claim of field-level performance. Our methods do not remove this bias; closing the lab-to-field gap would require cross-domain data and is left as future work.

\section{Data and Methods}\label{sec:methods}

\paragraph{Data.} PlantVillage \cite{mohanty2016using} contains 20{,}638 colour leaf images spanning 15 classes (pepper, potato, tomato; healthy and diseased). Class counts are strongly imbalanced: the largest class holds 3{,}208 samples (\textit{Tomato Yellow Leaf Curl Virus}) and the smallest only 152 (\textit{Potato healthy}); the full distribution is shown in Fig.~\ref{fig:class-dist} (Appendix). An exploratory pixel-level analysis (per-image mean/std, green-channel dominance, file-size $z$-scores) flagged no duplicated or corrupted samples; a small number of outliers were retained as within-class variation rather than artefacts.

\paragraph{Architecture.} A ResNet-50 backbone pretrained on ImageNet is followed by an embedding head \texttt{Linear(2048,1024) $\rightarrow$ BN $\rightarrow$ ReLU $\rightarrow$ Dropout(0.2) $\rightarrow$ Linear(1024,d)} with $\ell_2$-normalised output of dimension $d{=}512$. Classification at inference uses a prototype network: each class centroid is the mean embedding of its training samples and a query is assigned to the nearest centroid (cosine distance, equivalent to Euclidean on the unit sphere).

\paragraph{Training.} We minimise the triplet loss $\mathcal{L}_{\text{tri}}=\max(0,\,\|f_a-f_p\|_2-\|f_a-f_n\|_2+m)$ with margin $m$, mined online over batches drawn by an $M$-per-class sampler ($M{=}4$). The optimiser is Adam with a lower learning rate on the backbone than on the head; cross-validation uses 5 stratified folds, up to 50 epochs with early stopping on validation F1-macro (patience 7).

\paragraph{Hyperparameter optimisation.} Optuna with a MedianPruner ran 30 trials at 10-epoch fidelity on fold~0, jointly searching $d{\in}\{128,256,512\}$, $m{\in}[0.05,0.5]$, head learning rate in $[10^{-5},10^{-3}]$ (log), backbone-to-head LR ratio, weight decay, batch size $\in\{32,64,128\}$, and mining strategy $\in\{\text{batch-hard},\text{semi-hard},\text{multi-similarity}\}$. The selected configuration (Table~\ref{tab:hpo} in the Appendix) uses $d{=}512$, $m{=}0.178$, batch $32$, head LR $2.81{\times}10^{-4}$, backbone LR $8.65{\times}10^{-5}$, and \emph{multi-similarity} mining. To avoid HPO bias on fold~0, the final CV reports all five folds and fold~1 is used for downstream evaluation (confusion matrix, Grad-CAM, unknown detection).

\paragraph{Unknown detection (subtask~b).} We remove \textit{Tomato Bacterial spot} (2{,}127 images) from training, refit the embedding on the remaining 14 classes, and estimate a class-conditional Gaussian per class with a \emph{shared}, ridge-regularised covariance $\hat{\Sigma}+\lambda I$ \cite{lee2018simple}. The unknown score is $s(x)=\min_c (f(x)-\mu_c)^\top\hat{\Sigma}^{-1}(f(x)-\mu_c)$; a threshold $\tau$, chosen to maximise F1 on a held-out split, separates known from unknown.

\section{Results and Discussion}\label{sec:results}

\paragraph{Closed-set performance.} Across 5 stratified folds the mean F1-macro is $\mathbf{0.9985\pm0.0005}$ (Table~\ref{tab:cv}), well above the required 0.76 baseline. Per-class F1 is $1.0$ for eight classes and $\geq 0.996$ for all but one; the confusion matrix (Fig.~\ref{fig:cm-roc}, left) shows only single-digit misclassifications overall. This 0.998-level result is consistent with the CNN literature on PlantVillage \cite{ferentinos2018deep}; as noted in the introduction, it reflects in-distribution PlantVillage performance and inherits the dataset's lab-bias, and should not be taken as evidence of field-level generalisation.

\begin{table}[h]
\centering\small
\caption{5-fold stratified cross-validation, F1-macro.}\label{tab:cv}
\begin{tabular}{@{}lccccccc@{}}\toprule
Fold & 0 & 1 & 2 & 3 & 4 & Mean & Std \\\midrule
F1 & 0.9994 & 0.9980 & 0.9982 & 0.9984 & 0.9985 & \textbf{0.9985} & 0.0005 \\
\bottomrule
\end{tabular}
\end{table}

\paragraph{Error analysis.} The worst class is \emph{Potato healthy} (F1 = 0.9841), which is also the smallest class in the dataset (152 images, $0.74\%$): the prototype centroid is estimated from $\sim\!120$ training samples, which widens its confusion margin against \emph{Potato Late blight}. The next-worst classes, \emph{Tomato Target Spot} (0.9964) and \emph{Tomato Early blight} (0.9975), are visually similar brown-spotted tomato diseases and confuse symmetrically. Grad-CAM visualisations on fold~1 (Fig.~\ref{fig:gradcam}) confirm that the network attends to lesion regions rather than leaf silhouettes on correctly classified samples, while failure cases show broader leaf-surface activations.

\paragraph{Unknown detection.} With \emph{Tomato Bacterial spot} withheld, the Mahalanobis score cleanly separates known from unknown samples: the distance distributions (Fig.~\ref{fig:dist}) have means $4.67{\pm}1.57$ (known) vs.\ $21.74{\pm}9.40$ (unknown). AUROC is $\mathbf{0.9933}$ and PR-AUC $0.9899$; at the F1-optimal threshold $\tau{=}8.41$ the detector reaches F1 = $0.9505$ (precision $0.9489$, recall $0.9520$). The ROC curve (Fig.~\ref{fig:cm-roc}, right) is well above the diagonal across operating points, confirming that an open-set score built from the metric-learned embedding substantially outperforms a softmax-based one.

\begin{figure}[t]
\centering
\includegraphics[width=0.47\linewidth]{confusion_matrix.png}\hfill
\includegraphics[width=0.47\linewidth]{unknown_roc_curve.png}
\caption{Left: closed-set confusion matrix (fold~1). Right: ROC for unknown-class detection via Mahalanobis distance.}\label{fig:cm-roc}
\end{figure}

\paragraph{Conclusion.} Triplet-loss metric learning with a prototype classifier reaches $0.998$ F1-macro on PlantVillage and supports an effective Mahalanobis-based open-set detector (AUROC $0.993$) without retraining the representation. The remaining headroom is almost entirely in the rarest class, which additional samples or class-balanced sampling would address; a field-image evaluation is the natural next step.

\bibliography{sn-bibliography}

\begin{appendices}
\section*{Appendix A: Author Contributions}
\textbf{J.~Kraus} (corresponding): project lead; designed the end-to-end pipeline (data loading, triplet-loss training loop, prototype classifier, evaluation), ran HPO, produced the final cross-validation and unknown-detection results, and wrote the report. \textbf{R.~Schwarz}: exploratory data analysis, class-distribution and outlier-detection scripts, sample inspection of lab-bias artefacts. \textbf{D.~Charv\'at}: hyperparameter search-space design, Optuna integration and pruner configuration, training-curve visualisation. \textbf{M.~Kub\'ik}: subtask~(b) pipeline---held-out class setup, Mahalanobis scorer with shared covariance, ROC/PR evaluation and threshold selection.

\section*{Appendix B: Additional Figures and Tables}

\begin{table}[h]
\centering\small
\caption{Selected hyperparameters after 30 Optuna trials with MedianPruner (10-epoch trials on fold~0).}\label{tab:hpo}
\begin{tabular}{@{}ll@{}}\toprule
Parameter & Value \\\midrule
Embedding dim.\ $d$ & 512 \\
Triplet margin $m$ & 0.178 \\
Batch size & 32 \\
Head learning rate & $2.81\times 10^{-4}$ \\
Backbone learning rate & $8.65\times 10^{-5}$ \\
Weight decay & $5.37\times 10^{-5}$ \\
Mining strategy & multi-similarity \\
Optimiser & Adam (differential LR) \\
Sampler & $M$-per-class, $M{=}4$ \\
\bottomrule
\end{tabular}
\end{table}

\begin{figure}[h]\centering
\includegraphics[width=0.7\linewidth]{class_distribution.png}
\caption{PlantVillage class distribution (15 classes, 20{,}638 images). Smallest: \emph{Potato healthy} (152); largest: \emph{Tomato Yellow Leaf Curl Virus} (3{,}208).}\label{fig:class-dist}
\end{figure}

\begin{figure}[h]\centering
\includegraphics[width=0.48\linewidth]{training_curves_fold0.png}\hfill
\includegraphics[width=0.48\linewidth]{umap_embeddings.png}
\caption{Left: training and validation F1-macro per epoch (fold~0). Right: UMAP projection of the learned 512-d embeddings; clusters correspond to classes.}\label{fig:train-umap}
\end{figure}

\begin{figure}[h]\centering
\includegraphics[width=0.9\linewidth]{gradcam_samples.png}
\caption{Grad-CAM heatmaps on fold~1 samples. Activations concentrate on lesion regions for correctly classified images.}\label{fig:gradcam}
\end{figure}

\begin{figure}[h]\centering
\includegraphics[width=0.48\linewidth]{unknown_distance_histogram.png}\hfill
\includegraphics[width=0.48\linewidth]{unknown_confusion_matrix.png}
\caption{Left: Mahalanobis-distance histograms for known vs.\ unknown samples (\emph{Tomato Bacterial spot} held out). Right: known/unknown confusion matrix at the F1-optimal threshold $\tau{=}8.41$.}\label{fig:dist}
\end{figure}

\end{appendices}

\end{document}
